\documentclass[10pt,a4paper]{amsart}
\usepackage[margin=19mm]{geometry}
\usepackage{amsmath,amssymb,mathtools}
\usepackage[scr=rsfs,cal=euler]{mathalfa}
\usepackage{xfrac}
\usepackage[unicode,hidelinks,bookmarks=false]{hyperref}
\newcommand{\ie}{i.e.\ }
\newcommand{\cf}{cf.\ }
\newcommand{\resp}{respectively\ }
\newcommand{\Subsection}{\subsection}
\providecommand{\nobreakdash}{\nobreak\hbox{-}}
\setlength{\emergencystretch}{2em}
\multlinegap0pt
\usepackage{bm}
\usepackage{bbm}
\usepackage{dsfont}
\usepackage{xspace}
\usepackage{cancel}
\usepackage{mathtools}
\usepackage{amsthm}
\makeatletter
\@ifundefined{theorem}{\newtheorem{theorem}{Theorem}}{}
\@ifundefined{proposition}{\newtheorem{proposition}[theorem]{Proposition}}{}
\makeatother
\makeatletter
\expandafter\ifx\csname Forget\endcsname\relax
\newenvironment{Forget}{}{}
\fi
\newcommand{\boundary}[1]{\partial#1}
\DeclareMathOperator{\pr}{pr}
\DeclareMathOperator{\sing}{sing}
\makeatother
\title[Topological Complexity of symplectic CW-complexes]{Topological Complexity of symplectic CW-complexes}
\author{Luca Sandrock, Thomas Schick (Universit\"at G\"ottingen)}
\date{Source: arXiv:2505.19324v1 (CC BY 4.0)}
\begin{document}
\maketitle

\noindent\emph{Source and licence.} Topological Complexity of symplectic CW-complexes. Version arXiv:2505.19324v1 (CC BY 4.0), distributed under the Creative Commons Attribution 4.0 International licence, as recorded in the arXiv licence metadata for this exact version and shown on the abstract page https://arxiv.org/abs/2505.19324v1. Full rights notice: licenses/arxiv-2505.19324-CC-BY-4.0.txt.\par\smallskip

\noindent\textit{Editorial scope.} Counted unit: the Proposition whose source label is \texttt{prop:MV\_cocycle} in the article (source0.tex lines 399--414, source label \texttt{prop:MV\_cocycle}), delivered together with the article's own complete original proof of it (source0.tex lines 415--478, one \texttt{proof} environment of 64 lines). No mathematics is written, repaired or re-derived: statement and proof are the article's own text line for line. Required context retained uncounted: none. Every definition, display and label of those lines is retained unchanged. Omitted from this package and not redistributed: 1--398, 479--1282. Nothing inside a retained excerpt is omitted or altered: every retained body is delivered exactly as the article's lines are written, so no omission receipt of any kind accompanies this package. Citations: the retained excerpts carry no citation command at all, so no bibliography entry of the article is transcribed and no reference is redistributed. Reference adaptation: none was needed and none was made; every reference inside the delivered unit resolves inside it, so no unresolved reference or citation is produced. Printed slips kept literal: no printed word, symbol, hypothesis or number of the counted statement or of its proof was altered. Layout and class: the article's own class is \texttt{article} with packages including hyperref, amsrefs, amsmath, amsthm, amsfonts, latexsym, amscd, amssymb, enumerate; the wrapper is the lane's pinned \texttt{amsart} editorial wrapper at 10pt on A4, which restates the theorem-like environments the retained text needs and adds the article's own macro definitions that the retained text needs (\texttt{\textbackslash boundary}, \texttt{\textbackslash pr}, \texttt{\textbackslash sing}), with extra wrapper packages (bm, bbm, dsfont, xspace, cancel, mathtools). The delivered numbering is the unit's own. Assets: no retained span contains a figure environment, a graphics-inclusion command, a PDF-page inclusion, a table or a TikZ picture; no image file is copied into this package, the package declares no asset dependency and no image file is redistributed with it. Licence: version arXiv:2505.19324v1 of this article is distributed under the Creative Commons Attribution 4.0 International licence, as recorded in the arXiv licence metadata for this exact version and shown on the abstract page https://arxiv.org/abs/2505.19324v1. A full rights notice is delivered at licenses/arxiv-2505.19324-CC-BY-4.0.txt. Characters: every retained excerpt is pure ASCII, so the delivered engine, XeLaTeX with its default Latin Modern text font, reports no missing character.
\begin{proposition}\label{prop:MV_cocycle}
  Let $X$ be a connected CW-complex, $A$ an abelian coefficient group. Let $c\in
  C^2_{\sing}(X;A)$ be a singular cocycle (i.e.~$\boundary(c)=0$ for the
  singular cochain boundary map $\partial$). Set
  \begin{equation*}
    \overline c:= \pr_2^*c-
    \pr_1^*c \in C^2_{\sing}(X\times X; A)
  \end{equation*}
  and assume that there is a cochain $b_c\in
  C^1_{\sing}(PX;A)$ such that $\partial(b_c) = \pi^*\overline c$.

  Then $a_c:= r_1^*b_c-r_2^*b_c\in C^1_{\sing}(LX;A)$ is a cocycle and
  \begin{equation*}
    \delta_{MV}([a_c]) = \pi_2^*([\overline c]) \in H^2(P_2X;A).
  \end{equation*}
\end{proposition}

\begin{proof}
  First observe that $\pi\circ r_1=\pi\circ r_2$, hence by naturality
  \begin{equation*}
    \boundary(a_c) =r_1^*(\boundary b_c)-r_2^*(\boundary b_c) =
    r_1^*\pi^*\overline c - r_2^*\pi^*\overline c =0.
  \end{equation*}

  To compute the Mayer-Vietoris boundary map, we use its snake lemma
  definition applied to the Mayer-Vietoris short exact sequence of singular cochain groups
  {\small
    \begin{equation*}
  0 \to C^*_{\sing}(P_2X;A) \xrightarrow{i_1^*\oplus i_2^*}
  C^*_{\sing}(Z(r_1);A)\oplus C^*_{\sing}(Z(r_2);A) \xrightarrow{ \iota_1^*
    -\iota_2^*} C^*_{\sing}(LX;A) \to 0.
\end{equation*}
}

Specifically, given the cocycle $a_c\in C^1_{\sing}(LX;A)$, we need to find
  cochains $(\beta_1,\beta_2)\in C^1_{\sing}(Z(r_1);A)\oplus
  C^1_{\sing}(Z(r_2);A)$ such that $\iota_1^*\beta_1 - \iota_2^*\beta_2 =
  a_c$, where $\iota_k\colon LX\to Z(r_k)$ are the
  inclusion maps as end of the mapping cylinders.

  In our case, we let $\pr_k\colon Z(r_k)\to PX$ be the projection to the
  other end of the mapping cylinder (the standard homotopy equivalence) and
  use the cocycles
  \begin{equation*}
(\pr_1^*b_c,\pr_2^*b_c)\in C^1_{\sing}(Z(r_1);A)\oplus
  C^1_{\sing}(Z(r_2);A).
\end{equation*}
As we have $\pr_k\circ\iota_k = r_k$, by naturality we indeed get
\begin{equation*}
\iota_1^*\pr_1^*b_c- \iota_2^*\pr_2^*b_c= r_1^*b_c-r_2^*b_c = a_c.
\end{equation*}


We next have to compute the coboundary of $\pr_k^*b_c$. By naturality of the
coboundary map we obtain
  \begin{equation*}
    \boundary(\pr_k^*b_c) = \pr_k^*(\boundary b_c) = \pr_k^*\pi^*\overline c =
    i_k^*(\pi_2^* \overline c).
  \end{equation*}
Here, $i_k\colon Z(r_k)\to P_2X$ is the natural inclusion and the final
equality follows because upon composition with the endpoint projection to
$X\times X$, all relevant maps become equal: $\pi_2 \circ i_k = \pi\circ
\pr_k\colon Z(r_k)\to X\times X$. Now recall that $i_2^*\oplus i_2^*$ is the
inclusion map of the short exact sequence of singular cochain groups
{\small
  \begin{equation*}
  0 \to C^*_{\sing}(P_2X;A) \xrightarrow{i_1^*\oplus i_2^*}
  C^*_{\sing}(Z(r_1);A)\oplus C^*_{\sing}(Z(r_2);A) \xrightarrow{ \iota_1^*
    -\iota_2^*} C^*_{\sing}(LX;A) \to 0.
\end{equation*}
}
We just saw that
\begin{equation*}
  (i_1^*\oplus i_2^*) (\pi_2^* \overline c) =\boundary(\pr_1^*b_c,\pr_2^*b_c),
\end{equation*}
hence by the snake lemma indeed
\begin{equation*}
  \delta_{MV}([a_c]) =\pi_2^*[\overline c]
\end{equation*}
which is what we have to prove.
\end{proof}

\end{document}
